\section{Introduction}
\label{sec:introduction}

Group communication finds many use cases in blockchains, many of which play a vital role for the seamless
operation of their ecosystems.
Governance platforms~\cite{CIP-1694, DAO-White-Paper, eGov-DAO} broadcast voting-status updates, and discussions to 
Delegated Representatives (DReps), stakeholders and governance applications.
Stake pool operators need to inform their nodes-clients, about the results of their stake delegation,
by disseminating operational alerts, telemetry summaries or service-health information.
Decentralized Finance (DeFi)~\cite{Cardano-DeFi} applications need to disseminate market information, such as trading intentions,
price feeds, and liquidity alerts, through off-chain communication to support more educated trading decisions.

All the aforementioned use cases require a dedicated communication channel with several common characteristics.

First and foremost, to accommodate the decentralized nature of blockchain ecosystems, the communication system should itself be decentralized.
In contrast to centralized communication architectures, trust should not depend on a single trusted entity.
Receivers should be able to independently verify that a received event originates from an authorized sender.
To this end, the blockchain can serve as a shared registry of topic configurations and authorized senders, while digital signatures allow receivers to verify the authenticity of received events.

Secondly, the use-cases have in common that they require a single, or a few entities, sending an event to multiple.
In addition to that, the data sent, comes from observing the ecosystem, or the blockchain itself, thus there is no sensitive data.

Lastly, events should be disseminated promptly, and the nodes should never miss an event because of unexpected downtime, as they could be a crucial factor for a node to determine its decisions.

Taking the aforementioned requirements into consideration, we introduce \system, a decentralized Peer-to-Peer (P2P) publish/subscribe system that leverages the blockchain for topic administration and publisher authorization.
Event dissemination is carried out by the participating peers, while a separate set of dedicated replication servers, assumed to be operated by Stake Pool Operators (SPOs), provides persistent event storage and enables the recovery of missed events.
Our design and implementation focus on the Cardano ecosystem, thus, the terminology in the rest of the paper will be of the context of Cardano.

The contributions of this paper are threefold:
\begin{enumerate}[itemsep=6pt]
    \item We present the design and architecture of \system, detailing the key design decisions and explaining how they address the identified system requirements.
    \item We evaluate the proposed design through simulations, investigating its behavior and performance under different configurations.
    \item We implement an end-to-end client of \system, and and deploy it in a testbed to evaluate its performance under realistic workloads and failure scenarios.
\end{enumerate}

The rest of this paper is structured as follows.
\secref{sec:system-model} introduces the system model and the assumptions underlying our design.
\secref{sec:system-architecture} presents the architecture of \system and its components.
\secref{sec:evaluation} evaluates the system's performance through simulations and prototype-based experiments.
\secref{sec:related-work} surveys related work, and \secref{sec:conclusions} concludes the paper.
